For the modular quotient used below, let $N$ be the conductor of $A$.
Choose the finite surjective $\Q$-morphism
$\rho:X_0(N)\to A$ supplied by
\cite[Theorem C.13.6]{silverman}, and compose it with the finite
surjective forgetful morphism $u:X_1(N)\to X_0(N)$.
Put $f=\rho\circ u$. Pushforward of divisors preserves degree
and sends principal divisors to principal divisors by the function-field
norm \cite[Proposition II.3.6(c),(d)]{silverman}.
It therefore induces
\[
 q=f_*:J_1(N)=\operatorname{Pic}^0(X_1(N))
       \longrightarrow\operatorname{Pic}^0(A)\simeq A,
\]
where the last identification uses the identity point of $A$
\cite[Proposition III.3.4]{silverman}.
This is a homomorphism of abelian varieties over $\Q$: for the finite
flat map $f$, the norm on line bundles is
$\mathcal L\mapsto\det(f_*\mathcal L)\otimes
\det(f_*\mathcal O_{X_1(N)})^{-1}$;
this construction is defined over $\Q$, commutes with base change,
and induces the displayed map of Picard varieties.
To check surjectivity, work over $\overline{\Q}$ and fix any
$P_0\in X_1(N)(\overline{\Q})$. Given $a\in A(\overline{\Q})$,
choose $P$ with $f(P)=a+f(P_0)$. Then
$q([P]-[P_0])=a$. No rational point on $X_1(N)$ is required.
